% !PRINTED 602
\weberpage{602}

Thus a three-term cycle occurs in $Q$.

2. Let $\varkappa$ contain two three-term cycles $(1,\,2,\,3)(4,\,5,\,6)$. We take $\pi = (1,\,3,\,4)$ and obtain

\[
\lambda = (1,\,3,\,2)(4,\,6,\,5)(3,\,2,\,4)(1,\,5,\,6) = (1,\,2,\,5,\,3,\,4).
\]

This permutation $\lambda$, which is contained in $Q$, falls, however, under case 1.

3. Let $\varkappa$ contain a three-term and a two-term cycle $(1,\,2,\,3)(4,\,5)$. (That in $\varkappa$, if it belongs to $A$, a second cycle of even number of terms must also occur, is immaterial here). For $\pi = (1,\,2,\,4)$ we get

\[
\lambda = (1,\,3,\,2)(4,\,5)(2,\,4,\,3)(1,\,5) = (1,\,2,\,5,\,3,\,4),
\]

which again falls under case 1.

4. Let $\varkappa$ contain three transpositions, $(1,\,2)(3,\,4)(5,\,6)$. For $\pi = (1,\,3,\,5)$ it follows that

\[
\lambda = (1,\,2)(3,\,4)(5,\,6)(3,\,2)(5,\,4)(1,\,6) = (1,\,3,\,5)(2,\,6,\,4),
\]

which falls under case 2.

5. Let $\varkappa$ contain two transpositions and one unchanged element $(1,\,2)(3,\,4)(5)$. One sets $\pi = (1,\,2,\,5)$ and obtains

\[
\lambda = (1,\,2)(3,\,4)(5)(2,\,5)(3,\,4)(1) = (1,\,5,\,2).
\]

With this all cases are exhausted when $n > 4$. For $n = 4$ there still remains the one case in which $\varkappa$ consists of two transpositions, and it is just this that brings about the special behaviour at $n = 4$, by which the algebraic solution of the equation of degree $4$ is made possible\footnote{The first complete proof that the general equation of degree higher than the $4$th is not solvable by radicals is due to \gesperrt{Abel} (Crelle's Journal, Bd I, 1826). On the older attempts at a proof by \gesperrt{Ruffini} (1799 to 1806) and their relation to the \gesperrt{Abel} proof, compare the memoir of \gesperrt{Burkhardt}, ``Die Anfänge der Gruppentheorie und \gesperrt{Paolo Ruffini}'' (Abhandlungen zur Geschichte der Mathematik VI. Supplement zu Schlömilch's Zeitschrift. Leipzig 1892).}.

It follows further from this theorem that the symmetric group has no normal divisors other than itself, the alternating group and the unit group. For if $S$ is the symmetric group, $A$ the alternating group, and $Q$ a normal divisor of $S$, then the greatest common divisor $Q'$ of $A$ and $Q$ is a normal divisor of $A$, and is therefore equal to $A$ or equal to 1.
